\documentclass[mode=notes, palette=burgundy, lang=german, fontsize=11pt]{modernclassnotes}

\course{Theoretische Informatik \& Diskrete Mathematik}
\coursecode{CS 102}
\professor{Prof. Dr. Markus Weber}
\institution{Technische Universität München}
\term{Wintersemester 2026/27}
\title{Einführung in die Algorithmik}
\docsubtitle{Vorlesungsskript \& Übungsaufgaben}

\begin{document}

\maketitle
\tableofcontents
\newpage

\lecture[15. Oktober 2026]{1}{Komplexitätstheorie \& O-Notation}

\begin{definition}[Landau-Symbole]
Es seien $f, g: \mathbb{N} \to \mathbb{R}^+$ zwei Funktionen. Wir schreiben $f(n) = O(g(n))$, wenn positive Konstanten $c$ und $n_0$ existieren, sodass $f(n) \le c \cdot g(n)$ für alle $n \ge n_0$ gilt.
\end{definition}

\begin{theorem}[Hauptsatz der Laufzeitentwicklung (Master-Theorem)]
Seien $a \ge 1$ und $b > 1$ Konstanten, $f(n)$ eine Funktion und $T(n)$ definiert durch die Rekursionsgleichung $T(n) = a T(n/b) + f(n)$. Dann lässt sich das asymptotische Verhalten von $T(n)$ bestimmen.
\end{theorem}

\begin{proof}
Der Beweis erfolgt durch vollständige Induktion über die Rekursionsbäume der Aufteilungstiefe $\log_b n$.
\end{proof}

\begin{example}[Binäre Suche]
Die binäre Suche halbiert in jedem Schritt das Suchfeld, was zu einer Rekursion der Form $T(n) = T(n/2) + O(1)$ führt. Nach dem Master-Theorem gilt $T(n) = O(\log n)$.
\end{example}

\begin{exercise}[Laufzeitanalyse]
Bestimmen Sie die asymptotische Komplexität des Merge-Sort-Algorithmus mit der Rekursionsgleichung $T(n) = 2T(n/2) + O(n)$.
\end{exercise}

\begin{solution}
Durch Anwendung des Master-Theorems mit $a=2, b=2$ und $f(n)=O(n)$ erhalten wir $n^{\log_b a} = n^1 = n$. Da $f(n) = \Theta(n)$ ist, folgt $T(n) = \Theta(n \log n)$.
\end{solution}

\begin{takeaways}
\begin{itemize}
  \item Landau-Notation beschreibt das Wachstum von Funktionen im Unendlichen.
  \item Binäre Suche benötigt nur logarithmisches Wachstum $O(\log n)$.
  \item Master-Theorem ermöglicht die schnelle Analyse von Teile-und-Herrsche-Algorithmen.
\end{itemize}
\end{takeaways}

\end{document}
